\documentclass[10pt,a4paper]{amsart}
\usepackage{amsmath,amssymb,amsthm,graphicx}
\usepackage{verbatim}
\newenvironment{reference}{\comment}{\endcomment}
%\newenvironment{reference}{}{}
\newenvironment{slogan}{\comment}{\endcomment}
\newenvironment{history}{\comment}{\endcomment}

% For commutative diagrams we use Xy-pic
\usepackage[all]{xy}

% We use 2cell for 2-commutative diagrams.
\xyoption{2cell}
% Theorem environments.
%
\theoremstyle{plain}
\newtheorem{theorem}[subsection]{Theorem}
\newtheorem{proposition}[subsection]{Proposition}
\newtheorem{lemma}[subsection]{Lemma}

\theoremstyle{definition}
\newtheorem{definition}[subsection]{Definition}
\newtheorem{example}[subsection]{Example}
\newtheorem{exercise}[subsection]{Exercise}
\newtheorem{situation}[subsection]{Situation}

\theoremstyle{remark}
\newtheorem{remark}[subsection]{Remark}
\newtheorem{remarks}[subsection]{Remarks}

\numberwithin{equation}{subsection}

% Macros
%
\def\lim{\mathop{\mathrm{lim}}\nolimits}
\def\colim{\mathop{\mathrm{colim}}\nolimits}
\def\Spec{\mathop{\mathrm{Spec}}}
\def\Hom{\mathop{\mathrm{Hom}}\nolimits}
\def\Ext{\mathop{\mathrm{Ext}}\nolimits}
\def\SheafHom{\mathop{\mathcal{H}\!\mathit{om}}\nolimits}
\def\SheafExt{\mathop{\mathcal{E}\!\mathit{xt}}\nolimits}
\def\Sch{\mathit{Sch}}
\def\Mor{\mathop{\mathrm{Mor}}\nolimits}
\def\Ob{\mathop{\mathrm{Ob}}\nolimits}
\def\Sh{\mathop{\mathit{Sh}}\nolimits}
\def\NL{\mathop{N\!L}\nolimits}
\def\CH{\mathop{\mathrm{CH}}\nolimits}
\def\proetale{{pro\text{-}\acute{e}tale}}
\def\etale{{\acute{e}tale}}
\def\QCoh{\mathit{QCoh}}
\def\Ker{\mathop{\mathrm{Ker}}}
\def\Im{\mathop{\mathrm{Im}}}
\def\Coker{\mathop{\mathrm{Coker}}}
\def\Coim{\mathop{\mathrm{Coim}}}

% Boxtimes
%
\DeclareMathSymbol{\boxtimes}{\mathbin}{AMSa}{"02}

%
% Macros for moduli stacks/spaces
%
\def\QCohstack{\mathcal{QC}\!\mathit{oh}}
\def\Cohstack{\mathcal{C}\!\mathit{oh}}
\def\Spacesstack{\mathcal{S}\!\mathit{paces}}
\def\Quotfunctor{\mathrm{Quot}}
\def\Hilbfunctor{\mathrm{Hilb}}
\def\Curvesstack{\mathcal{C}\!\mathit{urves}}
\def\Polarizedstack{\mathcal{P}\!\mathit{olarized}}
\def\Complexesstack{\mathcal{C}\!\mathit{omplexes}}
% \Pic is the operator that assigns to X its picard group, usage \Pic(X)
% \Picardstack_{X/B} denotes the Picard stack of X over B
% \Picardfunctor_{X/B} denotes the Picard functor of X over B
\def\Pic{\mathop{\mathrm{Pic}}\nolimits}
\def\Picardstack{\mathcal{P}\!\mathit{ic}}
\def\Picardfunctor{\mathrm{Pic}}
\def\Deformationcategory{\mathcal{D}\!\mathit{ef}}
\usepackage{cleveref}
\crefname{originallemma}{original lemma}{original lemmas}
\Crefname{originallemma}{Original lemma}{Original lemmas}
\makeatletter
\newlabel{algebra-lemma-characterize-zero-local}{{00HN}{source}{}{sourcepointer-1}{}}
\newlabel{algebra-lemma-characterize-zero-local@cref}{{[originallemma][0][]00HN}{[1][0][]source}{}{}{}}
\makeatother
\crefname{originallemma}{original lemma}{original lemmas}
\Crefname{originallemma}{Original lemma}{Original lemmas}
\makeatletter
\newlabel{algebra-lemma-flat-going-down}{{00HS}{source}{}{sourcepointer-2}{}}
\newlabel{algebra-lemma-flat-going-down@cref}{{[originallemma][0][]00HS}{[1][0][]source}{}{}{}}
\makeatother
\crefname{originallemma}{original lemma}{original lemmas}
\Crefname{originallemma}{Original lemma}{Original lemmas}
\makeatletter
\newlabel{algebra-lemma-unique-prime-over-localize-below}{{00EA}{source}{}{sourcepointer-3}{}}
\newlabel{algebra-lemma-unique-prime-over-localize-below@cref}{{[originallemma][0][]00EA}{[1][0][]source}{}{}{}}
\makeatother
\title{A w-local morphism between w-local spectra which identifies local rings and is bijective on components is an isomorphism}
\author{The Stacks Project Authors}
\date{stacks-project:proetale.tex:local-snapshot-20261012; GNU FDL 1.2 or later}
\begin{document}
\maketitle
\noindent\textit{Editorial note.} All delivered mathematical text is the author's own, copied from the pinned Stacks Project chapter file proetale.tex; no statement or proof step was rewritten, reordered or removed.\par\medskip
\noindent\textit{Editorial note.} The counted proof is the author's complete argument: the bijection on the underlying spaces from the bijection on closed points and the decomposition of a w-local space into local rings at closed points, the comparison of the localizations through flatness and going down, and the closing local criterion for a ring map to be an isomorphism.\par\medskip
\noindent\textit{Editorial note.} The label lemma-w-local-morphism-equal-points-stalks-is-iso names the counted statement, and the proof environment immediately following it in the source is delivered verbatim as its proof.\par\medskip
\noindent\textit{Editorial note.} The counted statement is carried by the author's own lemma environment in the chapter source.\par\medskip
\noindent\textit{Editorial note.} Pre-open gate: python operations/agent-authoring/tools/precard\_gate.py --source proetale.tex --label lemma-w-local-morphism-equal-points-stalks-is-iso --start 1154 --end 1182 --local-id 28000 returned RESULT=OK (L28000 free, occupancy FREE, public interval clear). Candidate taken from sources/stacks-proetale/candidates-halfA.json.\par\medskip
\noindent\textit{Editorial note.} This card opens the second numbering segment L28000-L29999 (the first segment L3400-L3599 was closed by L3599); per LANE\_BRIEF section 4 already registered cards keep their numbers.\par\medskip
\noindent\textit{Editorial note.} The statement and proof are copied from the pinned chapter file; the cited inputs are kept as external pointers with their canonical Stacks tags.\par\medskip
\noindent\textit{Allowed standard prerequisites (editorial).}\begin{itemize}
\item Going down for flat ring maps. Algebra, the lemma that a flat ring map satisfies going down. Supplies the going down property of the flat map A -> B used to compare the localizations.
\item Uniqueness of the prime below for flat maps between localizations. Algebra, the lemma that for a flat map with going down the primes lying over a given prime are related as required, giving the identification of the localizations of B. Supplies the conclusion that for any prime q of B lying over p in A one has B\_q = B\_p.
\item Characterization of zero by vanishing of the localizations. Algebra, the lemma that a ring map is an isomorphism if it induces isomorphisms after localizing at every prime. Supplies the final step concluding A = B from the equality of all localizations.
\end{itemize}\medskip
\ifdefined\hypertarget\hypertarget{sourcepointer-1}{}\fi
\noindent\textit{Source reference.} Original lemma 00HN: sources/stacks-chow/tags.txt, source lines 649--649. Prerequisite pointer invoked at proof line(s) 1181; it names a standard fact of a sibling chapter that the counted argument uses, and every step of the counted proof built on it is retained verbatim in the delivered body.\par\medskip
\ifdefined\hypertarget\hypertarget{sourcepointer-2}{}\fi
\noindent\textit{Source reference.} Original lemma 00HS: sources/stacks-chow/tags.txt, source lines 654--654. Prerequisite pointer invoked at proof line(s) 1174; it names a standard fact of a sibling chapter that the counted argument uses, and every step of the counted proof built on it is retained verbatim in the delivered body.\par\medskip
\ifdefined\hypertarget\hypertarget{sourcepointer-3}{}\fi
\noindent\textit{Source reference.} Original lemma 00EA: sources/stacks-chow/tags.txt, source lines 527--527. Prerequisite pointer invoked at proof line(s) 1175; it names a standard fact of a sibling chapter that the counted argument uses, and every step of the counted proof built on it is retained verbatim in the delivered body.\par\medskip
\begin{lemma}
\label{lemma-w-local-morphism-equal-points-stalks-is-iso}
Let $A \to B$ be ring map such that
\begin{enumerate}
\item $A \to B$ identifies local rings,
\item the topological spaces $\Spec(B)$, $\Spec(A)$ are w-local,
\item $\Spec(B) \to \Spec(A)$ is w-local, and
\item $\pi_0(\Spec(B)) \to \pi_0(\Spec(A))$ is bijective.
\end{enumerate}
Then $A \to B$ is an isomorphism
\end{lemma}

\begin{proof}
Let $X_0 \subset X = \Spec(A)$ and $Y_0 \subset Y = \Spec(B)$ be the
sets of closed points. By assumption $Y_0$ maps into $X_0$ and
the induced map $Y_0 \to X_0$ is a bijection.
As a space $\Spec(A)$ is the disjoint union of the spectra
of the local rings of $A$ at closed points.
Similarly for $B$. Hence $X \to Y$ is a bijection.
Since $A \to B$ is flat we have going down
(Algebra, Lemma \ref{algebra-lemma-flat-going-down}).
Thus Algebra, Lemma \ref{algebra-lemma-unique-prime-over-localize-below}
shows for any prime $\mathfrak q \subset B$ lying over
$\mathfrak p \subset A$ we have $B_\mathfrak q = B_\mathfrak p$.
Since $B_\mathfrak q = A_\mathfrak p$ by assumption, we
see that $A_\mathfrak p = B_\mathfrak p$ for all primes $\mathfrak p$
of $A$. Thus $A = B$ by
Algebra, Lemma \ref{algebra-lemma-characterize-zero-local}.
\end{proof}

\end{document}
